\\newpage
\\section{问题分析}

\\subsection{问题一分析}
交通流时间序列具有复杂的动态特征：
\\begin{enumerate}
    \\item \\textbf{周期性}：日周期（早晚高峰）和周周期（工作日/周末）
    \\item \\textbf{趋势性}：长期增长或下降趋势
    \\item \\textbf{随机性}：天气、事件等随机因素影响
\\end{enumerate}
STL分解可以将序列分解为趋势项、季节项和残差项。

\\subsection{问题二分析}
Transformer模型的核心是Self-Attention机制，能够捕捉序列中的长距离依赖关系。对于交通流预测，我们需要：
\\begin{itemize}
    \\item 合理设计输入序列长度
    \\item 处理多变量输入（时间、日期、天气等）
    \\item 设计合适的损失函数
\\end{itemize}

\\subsection{问题三分析}
对比实验需要设计合理的评估指标（MAE、RMSE、MAPE）和 baseline 方法（ARIMA、LSTM），确保对比的公平性和科学性。
